% =====================================================================
% revisao_sistematica.tex — revisão sistemática (Parsifal). Entra logo abaixo de \section{Trabalhos Relacionados}.
% Marcadores para preencher: XXXX (números do Parsifal), as duas figuras e o número de artigos por extenso.
% =====================================================================

This systematic literature review was conducted using the Parsifal software. The first stage of the study involved defining the title and
objective of the research.

\subsection{Research Questions}

Subsequently, the following research questions were defined to guide which topics would be discussed in the selected final articles:

\begin{itemize}
\item[1)] quais são as aplicações atuais da integração de sistemas fuzzy com a incerteza das redes neurais profundas?
\item[2)] como a incerteza das redes neurais implementadas em sistemas fuzzy está sendo aplicada para enfrentar desafios específicos em
áreas de aprendizado de máquina?
\item[3)] quais são os avanços recentes na integração de redes neurais com a incerteza ou certeza, e como esses avanços podem ajudar na
classificação?
\item[4)] quais são os desafios e limitações que os modelos de redes neurais profundas enfrentam para reconhecer quando erram com alta
confiança, e como (ou se) sistemas fuzzy têm sido usados para estimar essa incerteza?
\end{itemize}

\subsection{Search String and Digital Libraries}

Then, the study's keywords were defined: ``Neural Networks'', ``deep learning'', ``machine learning'', ``fuzzy'', ``uncertainty'' and
``confidence''. The resulting search string was:

\begin{center}
\fbox{\parbox{0.94\columnwidth}{\footnotesize\ttfamily
("neural network*" OR "deep learning" OR "machine learning" OR CNN) AND fuzzy AND (uncertainty OR confidence OR certainty OR
hallucinat* OR overconfiden*)}}
\end{center}
% (ScienceDirect não aceita * e limita os conectores: se a string foi adaptada lá, citar em uma frase.)

The search string was used to filter articles in the following digital libraries: ACM Digital Library, IEEE Xplore, ScienceDirect and
Scopus. After the initial compilation of studies from the libraries, the articles were directly exported to Parsifal, with a total amount
of \textbf{XXXX} studies selected at the import stage.

\subsection{Selection Process}
\label{sec:revisao}

Following the databases search identifying \textbf{XXXX} articles, the strategy to filter the main studies is summarized in the
flowchart shown in Fig.~\ref{fig:fluxo}. Initially, duplicate studies were removed, prioritizing the counting of articles in the
databases where they were originally published. Subsequently, all articles went through the following exclusion criteria (CE):

\begin{figure}[t]
\centering
\IfFileExists{figuras/fluxograma_revisao.png}{\includegraphics[width=\columnwidth]{figuras/fluxograma_revisao.png}}{%
\fbox{\parbox[c][3.2cm][c]{0.92\columnwidth}{\centering\textcolor{gray}{Fluxograma de filtro de artigos\\
(salvar como figuras/fluxograma\_revisao.png)}}}}
\caption{Processo de seleção dos estudos.}
\label{fig:fluxo}
\end{figure}

\begin{list}{}{\setlength{\leftmargin}{2.9em}\setlength{\labelwidth}{2.5em}\setlength{\labelsep}{0.4em}}
\item[CE1:] estudos publicados antes de 2021;
\item[CE2:] estudos que não são artigos completos revisados por pares (resumos, pôsteres, capítulos, teses e revisões de literatura);
\item[CE3:] estudos em que o fuzzy não faz parte do método ou que não tratam incerteza ou confiança de redes neurais.
\end{list}

Na etapa seguinte, conforme apresentado no fluxograma, os estudos foram submetidos a uma análise de qualidade, detalhada na
Seção~\ref{sec:qualidade}, sendo mantidos aqueles que obtiveram pontuação máxima. Por fim, os \textbf{[número por extenso]} artigos
selecionados foram analisados para concluir o estudo, fornecendo uma resposta adequada às questões de pesquisa inicialmente propostas.

\subsection{Libraries Distribution from Exclusion Criteria}

Ao aplicar os critérios de exclusão e comparar os artigos iniciais com aqueles que permaneceram após a aplicação dos filtros CE1, CE2 e
CE3, foi elaborado o gráfico da Fig.~\ref{fig:bibliotecas}. À esquerda, ele apresenta o percentual de estudos importados de cada
biblioteca digital; à direita, o percentual de artigos remanescentes após a aplicação dos critérios de exclusão.

\begin{figure}[t]
\centering
\IfFileExists{figuras/bibliotecas_revisao.png}{\includegraphics[width=\columnwidth]{figuras/bibliotecas_revisao.png}}{%
\fbox{\parbox[c][3.2cm][c]{0.92\columnwidth}{\centering\textcolor{gray}{Gráficos de pizza: importados (esquerda) e remanescentes
após CE1--CE3 (direita)\\(salvar como figuras/bibliotecas\_revisao.png)}}}}
\caption{Distribuição dos estudos por biblioteca digital: importados (esquerda) e remanescentes após os critérios de exclusão (direita).}
\label{fig:bibliotecas}
\end{figure}

Inicialmente, observa-se que a biblioteca ACM Digital Library (em verde) apresenta o maior percentual de artigos importados. Entretanto,
após a aplicação dos critérios de exclusão, a biblioteca Scopus (em azul) passa a predominar, seguida pela IEEE Xplore (em amarelo).

\subsection{Quality Assessment}
\label{sec:qualidade}

% Perguntas sugeridas (nota 0 = não, 0,5 = em parte, 1 = sim). Ajuste ao que você cadastrou no Parsifal.
Os estudos remanescentes foram avaliados pelas seguintes questões de qualidade (QA), com nota 0 (não), 0,5 (em parte) ou 1 (sim):

\begin{list}{}{\setlength{\leftmargin}{2.9em}\setlength{\labelwidth}{2.5em}\setlength{\labelsep}{0.4em}}
\item[QA1:] o objetivo do estudo está claro?
\item[QA2:] o método está descrito com detalhe suficiente para ser reproduzido?
\item[QA3:] os experimentos usam datasets públicos?
\item[QA4:] o método é comparado com outras abordagens?
\item[QA5:] o estudo mede a incerteza ou o erro da rede (por exemplo, AUROC ou calibração), e não só a acurácia?
\end{list}
